\pdfbookmark[1]{Erläuterungen zu den Terminszeichnungen und den Beobachtungszahlen.}{erlaeuterungen-termine}
\begin{center}\large\textit{Erläuterungen zu den Terminszeichnungen und den Beobachtungszahlen.}\par
[\textit{Im Auszug.}]\end{center}
\begin{center}\small Resultate aus den Beobachtungen des magnetischen Vereins.\quad 1836.\quad V.\end{center}

Hier werden ---\quad---\quad---\quad---\quad--- die grafischen Darstellungen der Variationsbeobachtungen von sechs Terminen gegeben, insgesamt sechsundvierzig Kurven von vierzehn verschiedenen Beobachtungsorten: Berlin, Breda, Breslau, Catania, Freiberg, Göttingen, Den Haag, Leipzig, Mailand, Marburg, Messina, München, Palermo und Uppsala.
\begin{center}\small
---\enspace---\enspace---\enspace---\enspace---\enspace---\enspace---\enspace---\enspace---\enspace---\enspace---\enspace---\enspace---\enspace---\enspace---
\end{center}
Am 28. November 1835 und in der folgenden Nacht störte ein überaus heftiger Schirokko die Beobachtungen in Palermo so stark, dass sie einmal sogar für anderthalb Stunden unterbrochen werden mussten:
\SourceLeaf{575}Für viele Messreihen ließen sich nur einzelne, unzuverlässige Bestimmungen gewinnen. Daher ist zu vermuten, dass viele der sich ergebenden Schwankungen keine tatsächlich magnetischen Bewegungen waren. Wir wollten diese Kurve dennoch nicht ausschließen, da ihr letzter Teil vom Vormittag des 29. November, als sich der Sturm schon ziemlich gelegt hatte, eine ganz befriedigende Übereinstimmung mit den nördlichen Beobachtungsorten zeigt.

Bei dieser Gelegenheit sei noch bemerkt. dass nach allen sonstigen Erfahrungen selbst die heftigsten Sturmwinde keinerlei Wirkung auf die Magnetnadel haben, sofern der Beobachtungsraum und der Kasten dicht genug sind, um ihren unmittelbaren mechanischen Einfluss hinreichend abzuwehren. Im Göttinger magnetischen Observatorium wurde sehr oft beobachtet, dass sich die Nadel äußerst ruhig verhielt oder die Variation sehr gleichförmig fortschritt, während draußen der heftigste Sturm tobte. Wer jedoch nach solchen Erfahrungen gerade umgekehrt vermuten wollte, dass Stürme in der Atmosphäre die magnetischen Kräfte lähmen, würde durch den Verlauf des Januartermins 1836 widerlegt. Während dieses Termins herrschte in Göttingen und an mehreren anderen Beobachtungsorten ein sehr heftiger Sturm. Mehrere auswärtige Beobachter äußerten bei der Einsendung der Ergebnisse die Besorgnis, dass deshalb diesmal bei den ungemein starken Bewegungen wohl nur geringe Übereinstimmung zu finden sein werde. Gleichwohl stimmten bei diesem Termin, wie die Darstellung ---\quad---\quad---\quad--- zeigt, die Kurven von den verschiedenen Beobachtungsorten so vollkommen überein, dass man die Übereinstimmung bewundernswert nennen müsste, wäre sie nach so vielen Erfahrungen nicht etwas Gewohntes geworden. Nach mehreren Erfahrungen hier und anderswo haben Gewitter, selbst wenn sie nahe waren, ebenso wenig wie Stürme einen erkennbaren Einfluss auf die Magnetnadel gezeigt\expandafter\def\csname HistoricalNoteText575\endcsname{Natürlich ist hier nicht der Fall gemeint, in dem die atmosphärische Elektrizität mittels eines Zuleitungsdrahts durch einen die Nadel umgebenden Multiplikator zur Erde geführt wird.}\HistoricalNote{575}.

Ein im August 1836 eingegangenes Schreiben von Herrn von \textsc{Humboldt} enthielt die Nachricht, dass vom 10---18. August in Reykjavík auf Island der erfahrene französische Astronom Herr \textsc{Lottin} mit einem Apparat von \textsc{Gambey} ununterbrochen von Viertelstunde zu Viertelstunde die magnetische Variation beobachten werde, und den Wunsch, dass an einem oder einigen dieser Tage entsprechende Beobachtungen mit Magnetometern gemacht werden
\SourceLeaf{576}möchten. Daraufhin wurde für den 17---18. August ein außerordentlicher Termin angesetzt, und soweit es die Kürze der Zeit zuließ, wurden auswärtige Mitglieder unseres Vereins zur Teilnahme eingeladen. Dieser außerordentliche Termin wurde in Uppsala, Den Haag, Göttingen, Berlin, Leipzig und München genau auf die bei den ordentlichen Terminen eingeführte Art abgehalten. Und wenn die ---\quad---\quad---\quad---\quad--- grafisch dargestellten Beobachtungen recht interessante Bewegungen zeigen, so können wir nur bedauern, dass der ---\quad---\quad---\quad--- für die französischen Beobachtungen auf Island freigehaltene Platz leer bleiben musste, da wir über das Ergebnis dieser französischen Beobachtungen nichts erfahren konnten.

Der Septembertermin bietet eine Erfahrung, die hier etwas ausführlicher erwähnt werden soll, da sie das oben [S.\ 552 d.\ B.] Bemerkte auf lehrreiche Weise bestätigt. Im Protokoll der Marburger Beobachtungen, die diesmal in Abwesenheit von Herrn Professor \textsc{Gerling} ohne dessen persönliche Teilnahme ausgeführt worden waren, fanden sich für $12^{\mathrm h}5'$ nur ganz unregelmäßig verlaufende Zahlen, die überhaupt kein Ergebnis liefern. Für $12^{\mathrm h}10'$ erscheint auf einmal eine um $30{,}54$ Skalenteile größere Zahl als für $12^{\mathrm h}0'$. ---\quad---\quad---\quad---\quad---\quad---\quad---\quad---\quad---\quad---\quad---\quad--- Das weckte die Vermutung, dass um $12^{\mathrm h}5'$ eine Spinne die freie Bewegung der Nadel gehemmt habe, indem sie einen Faden angeknüpft hatte. Diese Vermutung wurde noch wahrscheinlicher, weil von $12^{\mathrm h}10'$ bis zum Ende die Bewegungen der Nadel zwar denen, die sich aus den Beobachtungen an anderen Orten ergaben, ganz ähnlich waren, aber verhältnismäßig viel kleiner ausfielen, als man nach den Erfahrungen bei anderen Terminen hätte erwarten müssen. Herr Professor \textsc{Gerling} wurde deshalb gebeten, nach seiner Rückkehr nach Marburg den Apparat genau zu untersuchen. Das Ergebnis sei hier noch aus einem Schreiben von Herrn Professor \textsc{Gerling} vom 12. November mitgeteilt.

\Needspace{9\baselineskip}
Die Untersuchung wurde am 5. November vorgenommen. Seit dem Septembertermin hatte bis dahin niemand das Beobachtungszimmer wieder betreten. Zuerst wurde die Stellung der Nadel bestimmt. Man fand
\begin{center}
\begin{tabular}{rrc l}
$3^{\mathrm h}$&$33'$&.\quad.\quad.\quad.\quad.\quad.&$445{,}63$ Skalenteile\\
&35&.\quad.\quad.\quad.\quad.\quad.&$445{,}73$\\
&37&.\quad.\quad.\quad.\quad.\quad.&$445{,}71$
\end{tabular}
\end{center}

\SourceLeaf{577}Darauf wurde die Nadel mithilfe des sogenannten Beruhigungsstabs in mäßige Schwingungen versetzt. Daraus ergab sich eine Schwingungsdauer von 17 Sekunden, neun Sekunden weniger als die sonst bekannte Schwingungsdauer. Als darauf der Deckel des Kastens vorsichtig abgehoben wurde, bemerkte man an seiner Unterseite eine sehr kleine lebende Spinne. Man glaubte auch, einen daran hängenden, wenn auch kaum wahrnehmbaren Faden zu sehen. Ferner fand man im Kasten eine Anzahl kleiner schwarzer, punktartiger Körper, die sich unter dem Mikroskop als tote Mücken erwiesen. Schließlich fand man in einer Ecke des Kastens ein regelrechtes, unversehrtes Gespinst, das so fein war, dass es ohne den Widerschein der Lichter kaum erkennbar gewesen wäre. Nach allen Umständen konnte man nur annehmen, dass sich die Spinne schon seit längerer Zeit im Kasten aufgehalten hatte.

Nachdem man dann noch mit dem Finger an allen Seiten des Magnetstabs entlanggefahren war, ergaben neue Beobachtungen der Schwingungsdauer wieder genau den alten Wert von 26 Sekunden. Auch fand man wieder eine erheblich niedrigere Stellung als zuvor, nämlich
\begin{center}
\begin{tabular}{rrc l}
$4^{\mathrm h}$&$43'$&.\quad.\quad.\quad.\quad.\quad.&$431{,}15$ Skalenteile\\
&45&.\quad.\quad.\quad.\quad.\quad.&$431{,}46$\\
&47&.\quad.\quad.\quad.\quad.\quad.&$431{,}12$
\end{tabular}
\end{center}
Natürlich lässt sich aus diesen Beobachtungen der Stellung nicht genau bestimmen, um wie viel das nun beseitigte Hindernis die Stellung verfälscht hatte. Denn die etwaige Änderung der Deklination während der mehr als eine Stunde dauernden Zwischenzeit blieb unbekannt.

Bei dieser Gelegenheit sei noch ein zweiter Vorfall ähnlicher Art erwähnt. Die Schwingungsdauer des Magnetstabs in Breslau, die im März 1836 beinahe 32 Sekunden betrug, war von da bis zum November ganz allmählich größer geworden und hatte in dieser Zeit insgesamt um etwa $0''4$ zugenommen. Das ist ganz in Ordnung, da alle Magnetstäbe im Laufe der Zeit etwas von ihrer Kraft verlieren, wenn auch in sehr unterschiedlichem Maß, je nach Härtegrad des Stahls und anderen Umständen. Vom November 1836 bis Januar 1837\expandafter\def\csname HistoricalNoteText577\endcsname{Vermutlich wurde die Schwingungsdauer in der \textit{Zwischenzeit} nicht bestimmt.}\HistoricalNote{577} war dagegen wieder eine
\SourceLeaf{578}Abnahme der Schwingungsdauer um $1''27$ eingetreten. Und Herr Professor von \textsc{Boguslawski}, der mir diesen auffallenden Umstand in einem Schreiben vom 5. März mitteilte, schien geneigt, ihn zum Teil einer größeren Intensität des Erdmagnetismus zuzuschreiben. Mir schien jedoch außer Zweifel, dass der Grund für dieses Phänomen in der nächsten Umgebung des Magnetstabs zu suchen war, wahrscheinlich darin, dass er sich nicht ganz frei bewegen konnte, obwohl von einem Spinnenfaden nach den Marburger Erfahrungen eher eine erheblich stärkere Wirkung zu erwarten gewesen wäre. Auch Herr von \textsc{Boguslawski} fand diese Vermutung bestätigt. Unter dem 21. März antwortete er: „Die Ursache der Änderung der Schwingungsdauer haben Sie richtig erraten. Der Kasten war zufällig etwas zur Seite geschoben worden, sodass der Rand des kleinen Lochs, durch das der Faden von oben eintritt, diesem nahe gekommen war, ihn aber keineswegs berührte. Dennoch müssen einige feine Fasern bis an den Rand gereicht haben, denn seit der Faden wieder durch die Mitte des Lochs geht, ist auch die Schwingungsdauer wieder annähernd dieselbe wie früher.“

Zu den Bewegungen selbst, die hier aus sechs Terminen dargestellt werden, seien noch einige Bemerkungen hinzugefügt.

In den drei Sommerterminen ---\quad---\quad---\quad---\quad---\quad---\quad--- sieht man durch alle großen Anomalien hindurch doch auch die regelmäßige tägliche Bewegung, insofern die Kurven in den Nachmittagsstunden ansteigen und in den folgenden Vormittagsstunden abfallen. In den drei Winterterminen dagegen ---\quad---\quad---\quad---\quad---\quad---\quad--- ist davon kaum noch etwas zu erkennen. Nach all unseren Erfahrungen ist es tatsächlich sehr gewöhnlich, dass das Regelmäßige vom Unregelmäßigen überragt wird oder ganz darin untergeht. In den Jahren 1834 und 1835 gab es jedoch auch einige Termine, an denen keine größeren Anomalien den regelmäßigen Verlauf verdunkelten, während kleine nie fehlten.

Was aber die anomalen Bewegungen so bemerkenswert macht, ist die außerordentlich große Übereinstimmung an verschiedenen, ja meist an sämtlichen Orten. Sie reicht gewöhnlich bis in die feinsten Nuancen, nur die Größenverhältnisse sind unterschiedlich. Es wäre ganz unnötig, diese Übereinstimmung im Einzelnen nachzuweisen. Der Anblick unserer sechs Termindarstellungen spricht schon hinreichend für sich.

\SourceLeaf{579}Vorerst kann es noch gar nicht unsere Aufgabe sein, diese rätselhafte Hieroglyphenschrift der Natur zu entziffern. Zunächst müssen wir uns darauf beschränken, Abschriften dessen zu sammeln, was sich darbietet, und sie immer zuverlässiger, getreuer und vielfältiger zu machen. Wenn reichlich Material vorliegt, wird es, wie wir zuversichtlich hoffen dürfen, eines Tages auch an der Entzifferung nicht fehlen. Inzwischen sei es erlaubt, einige Bemerkungen hinzuzufügen, die zu einer zutreffenderen Beurteilung beitragen können.

Zunächst darf man nicht vergessen, dass alle solche Anomalien vergleichsweise nur geringe Veränderungen oder Zusätze zur großen erdmagnetischen Kraft sind (oder, genau gesprochen, zu deren horizontaler Komponente), dass wir zwischen diesen Anomalien und jener Kraft sorgfältig unterscheiden müssen und dass uns nach dem bisherigen Stand nichts dazu zwingt, beide auf gleiche oder gleichartige Ursachen zurückzuführen. Man mag es immerhin für wahrscheinlich halten --- was wir ganz dahingestellt lassen ---, dass jene Anomalien Wirkungen elektrischer Ströme oder Ausgleichsvorgänge sein können, vielleicht weit außerhalb der Atmosphäre. Deshalb braucht man die ältere Vorstellung noch nicht aufzugeben, dass die Hauptkraft im festen Teil des Erdkörpers selbst ihren Sitz hat oder vielmehr die Gesamtwirkung aller magnetisierten Teile des Erdkörpers ist. Wäre das Erdinnere, wie einige Naturforscher meinen, noch flüssig, so böten die fortschreitende Verfestigung und die daraus folgende Verdickung der festen Erdkruste die natürlichste Erklärung für die Säkularänderungen der magnetischen Kraft.

Wir verlassen jedoch lieber den unsicheren Boden der Hypothesen und kehren zu den Tatsachen zurück. Bei Weitem die meisten Anomalien finden wir an den südlicheren Beobachtungsorten kleiner und an den nördlicheren größer. So beträgt zum Beispiel der bemerkenswerte Anstieg am 30. Januar 1836 zwischen $9^{\mathrm h}25'$ und $9^{\mathrm h}40'$, auf Bogenteile umgerechnet, in Catania 6, in Mailand 12, in München $13\frac12$, in Leipzig 16, in Marburg 20, in Göttingen 26 und in Den Haag 29 Bogenminuten. Von diesem Unterschied ist allerdings etwas abzuziehen, weil an den nördlicheren Punkten, wo die horizontale Komponente der erdmagnetischen Kraft selbst eine geringere Intensität hat als an den südlicheren, gleiche Störkräfte eine stärkere Wirkung haben müssen als an den südlicheren Punkten. Der Unterschied der Intensitäten von Den Haag bis Catania ist jedoch im Vergleich zu den beobachteten Unterschieden nur gering. Es steht also fest, dass
\SourceLeaf{580}die Stärke der damaligen Störkraft umso geringer war, je weiter nach Süden wir ihre Wirkung verfolgen. So wenig wir über das Wesen solcher Störkräfte wissen, können wir doch nicht daran zweifeln, dass jede von ihnen ihre Quelle an einem bestimmten Ort im Raum haben muss. So wie wir den Sitz der Kraft, die die erwähnten Erscheinungen hervorgebracht hat, notwendig nördlich oder nordwestlich der Beobachtungsorte annehmen müssen (ohne bei so wenigen Daten eine nähere Bestimmung zu wagen), so scheinen überhaupt die nördlichsten Gegenden das Hauptgebiet zu sein, von dem die meisten und stärksten Wirkungen ausgehen. Dies gilt, soweit man derartige Schlüsse schon auf Erfahrungen aus einem im Verhältnis zur ganzen Erdoberfläche doch nur kleinen Gebiet stützen darf.

Betrachtet man das bisher vorliegende Material genauer, so zeigen sich bei den verschiedenen aufeinanderfolgenden Bewegungen beträchtliche Unterschiede in ihren Größenverhältnissen an verschiedenen Orten, selbst wenn die Ähnlichkeit sonst ganz unverkennbar ist. So ist zum Beispiel von zwei kurz aufeinanderfolgenden Ausschlägen an einem Ort oft der erste größer, an einem anderen Ort umgekehrt. Wir müssen daher annehmen, dass am selben Tag und in derselben Stunde viele Kräfte zugleich wirksam sind, die vielleicht ganz voneinander unabhängig sind und sehr verschiedene Sitze haben. Ihre Wirkungen können sich an den verschiedenen Beobachtungsorten je nach Lage und Entfernung in sehr unterschiedlichen Verhältnissen vermengen oder ineinandergreifen, wenn die eine zu wirken beginnt, bevor die andere aufgehört hat. Die Verwicklungen aufzulösen, die so in den Erscheinungen an jedem einzelnen Ort entstehen, wird zweifellos sehr schwer sein. Gleichwohl dürfen wir zuversichtlich hoffen, dass diese Schwierigkeiten nicht immer unüberwindlich bleiben werden, wenn sich die Teilnahme an den gleichzeitigen Beobachtungen noch viel weiter verbreitet haben wird. Es wird der Triumph der Wissenschaft sein, wenn es eines Tages gelingt, das bunte Gewirr der Erscheinungen zu ordnen, die \textit{einzelnen} Kräfte auseinanderzulegen, deren zusammengesetztes Ergebnis sie sind, und für jede ihren Sitz und ihr Maß nachzuweisen.

Nicht ganz selten findet man auch an einzelnen Orten einen kleinen Ausschlag, zu dem sich an den übrigen Orten kein Gegenstück erkennen lässt. Es wäre jedoch zu gewagt, dergleichen sofort als bloß lokale magnetische
\SourceLeaf{581}Einwirkung zu erklären. Bei so vielen Zahlen kann tatsächlich gelegentlich ein Irrtum vorgekommen sein. Uns sind öfter solche Fälle begegnet, in denen ein Blick in die Originalbeobachtungen, sofern wir sie zur Hand hatten, einen Rechenfehler bei der Reduktion oder einen offensichtlichen Schreibfehler erkennen ließ. In ähnlichen Fällen, in denen wir nur den Auszug aus den Beobachtungen vorliegen hatten, hatte es zuweilen denselben Erfolg, den Einsender auf einen solchen Verdacht hinzuweisen. Da es jedoch nicht praktikabel ist, alle derartigen Fälle immer erst im Briefwechsel zu erörtern, können diejenigen Teilnehmer, die die Originalbeobachtungen selbst nicht einsenden, solche Stellen in den Kurven ihrer Beobachtungen zum Anlass nehmen, in den Originalen nachzusehen, etwa für Leipzig am 26. November 1836 um $6^{\mathrm h}15'$ Göttinger Zeit. Berichtigungen, die sich auf diese Weise ergeben, sollen dann nachträglich mitgeteilt werden. Völlige Gewissheit hat man bei solchen Stellen, die nur auf einem einzigen Satz von Beobachtungen beruhen, aber selbst dann noch nicht, wenn die Originalaufzeichnungen keinen Fehler eindeutig nachweisen. Denn auch einem nicht ganz ungeübten Beobachter kann es einmal passieren, dass er im selben Satz wiederholt falsche Zehnerstellen der Skalenwerte aufschreibt. Mit einer solchen, freilich etwas gewagten Vermutung ließe sich etwa die oben erwähnte Zahl von $11.69$ auf $6.69$ bringen und so in die übrigen Werte einfügen.

Für lokal im \textit{engsten} Sinn dürfte man einen solchen isolierten Ausschlag übrigens auch dann noch nicht halten, wenn die Tatsache selbst keinem Zweifel mehr unterliegt. Wie die Quelle jeder Anomalie irgendwo ihren Sitz haben muss, so kann der Sitz der einen oder anderen auch einmal nahe bei einem der Beobachtungsorte selbst liegen. Ist eine solche Kraft an sich nur schwach, so kann ihre Wirkung am nächsten Ort gerade wegen der Nähe augenfällig sein und an allen übrigen Orten, \textit{an denen beobachtet wird}, verschwinden (das heißt für uns nicht mehr erkennbar sein), eben weil diese schon \textit{zu} weit entfernt sind. Zumindest bisher scheint es daher gar keinen Grund zu geben, zwischen den Anomalien andere als quantitative Unterschiede anzunehmen. Zugleich folgt daraus, dass es in manchen Fällen sehr nützlich sein kann, wenn zwei oder mehrere Beobachtungsorte nur mäßig weit voneinander entfernt liegen. So wäre es zum Beispiel sehr wünschenswert gewesen, wenn in Augsburg (wo man jetzt regelmäßig an den Terminsbeobachtungen teilnimmt) bereits der Septembertermin
\SourceLeaf{582}1836 beobachtet worden wäre. Sehr wahrscheinlich hätte man dann die Bewegung um $2^{\mathrm h}10'$, die zwar an den meisten Orten durchschimmert, aber in München auffallend größer ist, mit mehr Sicherheit beurteilen können, als es jetzt möglich ist.
\begin{center}\small
---\enspace---\enspace---\enspace---\enspace---\enspace---\enspace---\enspace---\enspace---\enspace---\enspace---\enspace---\enspace---\enspace---\enspace---
\end{center}

